%%%%% Section 1 Introduction %%%%%
\section{Introduction}\label{sec:introduction}

A visitor comes to the reference desk of a music library.

\begin{quote}
I am looking for something like Miles Davis---that atmosphere I cannot
quite name.
\end{quote}

The librarian asks for something more specific. The visitor searches for
words, falls silent for a while, and finally offers:

\begin{quote}
Quiet, but with feeling in it; impressionistic; colours that seem to melt
into one another\ldots{} I cannot say it well.
\end{quote}

A visitor who gets this far is doing well; most never move past the first
sentence. The librarian understands the request---or feels she does. But
consider what the catalogue gives her to work with. From ``Miles Davis'' it
retrieves a genre (jazz), an instrument (trumpet), a period (the 1950s and
60s), a country (the United States). A search on these facets returns hard
bop, fusion, and a quantity of other trumpeters, most of it not what the
visitor wants. Some weeks later, by another route---browsing a shelf, or a
friend's recommendation---he comes upon Debussy's piano music, and thinks:
\emph{so this is what I was looking for.}

This scene is our point of departure, and we take it apart before building
on it. First, everything the catalogue records has changed. Jazz has become
classical, trumpet piano, the 1950s the 1900s, America France; a performer
has even become a composer. Not one catalogued attribute survived the
passage---and yet the visitor's criterion was met. Second, the affinity is
not a private fancy. The kinship between the music around Miles Davis and
impressionist harmony has been remarked many times: the modal language of
\emph{Kind of Blue}, Bill Evans's pianism, and Gil Evans's orchestration
have all been discussed in relation to Debussy and
Ravel~\cite{TODO-musicology}. What the visitor sensed lies \emph{in the
works}; it is simply written in none of the catalogue's fields. Third, what
the visitor discovered was not a piano piece. He did not say ``what a fine
piece''; he said ``this is what \emph{I} was looking for.'' The object of
his discovery was his own criterion---something he could not have stated
before the encounter. Fourth, he had indicated his need in two ways: by
pointing at a work (``like Miles Davis''), and by heaping up partially apt
words (``quiet, with feeling, melting colours''). Neither is a subject
term, and the catalogue has a field for neither.

Wilson distinguished two powers at which bibliographic control can
aim~\cite{Wilson1968}: \emph{descriptive} power, the power to find
everything that bears on a topic, and \emph{exploitative} power, the power
to put recorded knowledge to one's own best use. What the cataloguing
tradition built to a high finish---from Cutter's objects through facet
analysis---is the former; the latter has been left, for the most part, to
the skill of readers and librarians. Our visitor is asking for the latter:
not what a work is \emph{about}, but how it \emph{feels}.

Two kinds of attribute must therefore be told apart, and the pair will
anchor everything that follows. \emph{Denotative} attributes can be
expressed as sets of mutually exclusive labels and implemented as facets:
genre, period, creator, instrument, country. \emph{Connotative} attributes
cannot. They resist a single name, come in degrees, are many-dimensional,
interfere with one another, and take form in the encounter between work and
listener. What Miles Davis and Debussy share is of this second kind. Cross
the pair with the familiar division between known-item search and
exploratory browsing, and four quadrants appear
(Table~\ref{tab:quadrants}). Faceted navigation is exploratory yet entirely
denotative, so the customary opposition of ``facets versus exploration''
misses the point: what facets fail to carry is not exploration but the
unnamed. And the quadrant of connotative search---retrieving a work by a
quality no one can name---is nearly empty. Why it is empty will become
clear in Section~\ref{subsec:third}.

\begin{table}[t]
\centering
\caption{Two kinds of attribute crossed with two modes of seeking.}
\label{tab:quadrants}
\begin{tabular}{lll}
\toprule
 & Known-item search & Exploratory browsing \\
\midrule
Denotative  & bibliographic search, OPAC & faceted navigation \\
Connotative & (nearly empty) & readers' advisory, Book House \\
\bottomrule
\end{tabular}
\end{table}

This problem is not new to library and information science. It has been
worked on from three directions, by three lineages that rarely cite one
another. How ``quiet, with feeling, melting colours'' might be written into
a bibliographic record is a question of \emph{representation}
(Section~\ref{subsec:first}). That the visitor could not say what he sought
until he had found it is a question about \emph{the seeker}
(Section~\ref{subsec:second}). How a system might receive ``like Miles
Davis'' is a question about \emph{the means of access}
(Section~\ref{subsec:third}). Section~\ref{sec:lineages} follows each
lineage to its endpoint; laid over one another, the three endpoints
determine what is required almost uniquely. It is not another interface. It
is a \emph{representation}---something written into records, outliving its
maker, shared across libraries---for qualities that have no name; for no
interface can surface what the representation beneath it does not already
encode~\cite{Blair1985289,Furnas1987964}. Section~\ref{sec:framework}
states the proposal: a coordinate system in which a work's nameable
properties and its unnamed qualities occupy separate coordinates, the
unnamed ones serving as \emph{headings without names}---and there
serendipity, too, acquires a precise sense: the seeker's recognition that
one of their own unnamed targets is expressed by such a coordinate.
Section~\ref{sec:mathematics} supplies the mathematics on which the
proposal leans---the metric under which the coordinate system is unique, so
that whoever computes it obtains the same one; this, we will argue, is the
condition for coordinates to serve as authority data.
Section~\ref{sec:computation} shows that the coordinates can be learned
from a collection and its catalogue metadata,
Section~\ref{sec:demonstration} demonstrates this on real data,
Section~\ref{sec:discussion} discusses what changes for libraries, and
Section~\ref{sec:conclusion} concludes.
